% Generated draft: requires mathematical and source review.
\documentclass[11pt]{article}
\usepackage[margin=25mm]{geometry}
\usepackage{fontspec}
\setmainfont{Latin Modern Roman}
\usepackage{amsmath,amssymb,mathtools,unicode-math}
\setmathfont{Latin Modern Math}
\usepackage{enumitem,needspace,xurl,hyperref}
\setlength{\parindent}{0pt}
\setlength{\parskip}{5pt}
\setlength{\emergencystretch}{3em}
\newcommand{\N}{\mathbb{N}}
\newcommand{\Z}{\mathbb{Z}}
\newcommand{\Q}{\mathbb{Q}}
\newcommand{\R}{\mathbb{R}}
\newcommand{\C}{\mathbb{C}}
\newcommand{\F}{\mathbb{F}}
\newcommand{\A}{\mathbb{A}}
\newcommand{\PP}{\mathbb{P}}
\newcommand{\E}{\mathbb{E}}
\newcommand{\eps}{\varepsilon}
\newcommand{\dd}{\,\mathrm{d}}
\begin{document}
\section*{UnsolvedMath ID: 30006310}
\textbf{Source problem code:} catalogue entry
\subsection{Definitions and mathematical statement}
Let \(k\) denote a ground field, whose characteristic and algebraic closure assumptions are not specified in the supplied record. A toric variety \(X\) has an algebraic torus \(T\) as a dense open subset, with the multiplication action extending to \(X\); the precise convention concerning splitness and normality must be fixed for a definitive formulation. A toric vector bundle \(\mathcal E\) on \(X\) is a finite-rank locally free sheaf equipped with a lift of the \(T\)-action that acts linearly on the fibers of its associated vector bundle. For positive rank, the quotient convention for projectivization is
\[
\pi\colon Y=\PP_X(\mathcal E)=\operatorname{Proj}_X\!\bigl(\operatorname{Sym}_{\mathcal O_X}\mathcal E\bigr)\longrightarrow X,
\]
which parametrizes one-dimensional quotients of the fibers of \(\mathcal E\). The alternative convention parametrizes one-dimensional subspaces and replaces \(\mathcal E\) by \(\mathcal E^\vee\) in this formula. The source does not select a convention.

\textbf{Research question, with unresolved scope.} For projectivizations \(Y\) of toric vector bundles, characterize the structures, invariants, and vanishing properties that can be compared with those of smooth projective toric varieties, and determine the hypotheses under which each proposed comparison holds.

This is a research program rather than a fully specified conjecture. The record does not identify particular structures, invariants, cohomology groups, or vanishing assertions. It also does not specify whether \(X\) is smooth, complete, or projective, the allowed ranks of \(\mathcal E\), or whether \(Y\) is required to be smooth or projective. The original wording uses ``a (toric) vector bundle,'' leaving the possible inclusion of bundles without a torus-equivariant structure unresolved. Consequently, no universal comparison theorem, precise vanishing statement, or assertion that \(Y\) itself is toric can be extracted without additional choices.
\subsection{Short English statement}
The question asks which features of smooth projective toric varieties also occur for projectivized toric vector bundles, and under which hypotheses. The supplied wording leaves the comparison properties and the class of bundles and base varieties unspecified, so it does not determine a single precise conjecture.
\subsection{Sources}
\begin{itemize}
\item Supplied local UnsolvedMath registry record, permanent ID \(30006310\), source problem code \texttt{catalogue entry}, entitled ``Projectivized Toric Vector Bundles as Toric Varieties,'' in the Oberwolfach Reports Open Problems collection. The record labels the problem open and its statement \texttt{corrected\_verified}; these are supplied curation labels, not independent mathematical verification. Its verification note states that only formula delimiters, notation, grammar, and typography were repaired, with mathematical content and scope preserved.
\item Supplied source citation: \emph{Toric Geometry} (2025), Oberwolfach Reports, DOI 10.4171/owr/2025/19, pp.\ 883--940. \url{https://doi.org/10.4171/owr/2025/19}. The supplied background describes these reports as documenting workshops and their open-problem discussions. No entry-specific page is supplied.
\item The original extract attributes the surrounding contribution to Gregory G. Smith, joint work with Michael Perlman, and appends the phrase ``TBD: Toric Bundles Duh!'' The dedicated proposer field is empty; the supplied year is 2025. The precise role of the appended phrase is unspecified.
\item Supplied literature source: Bakhtary, \emph{Canonical Projectivization of Smooth Complete Toric Varieties}. \url{https://arxiv.org/abs/2605.25013}.
\item Supplied literature source: Gonzalez-Anaya et al., \emph{Dynamics of Projectivized Toric Vector Bundles}. \url{https://arxiv.org/abs/2510.26757}.
\item Literature assessment dated exactly \texttt{2026-08-22}: supplied dated curation, not independently verified. It labels the question open and describes existing work as addressing special bundles and related projectivization or endomorphism problems rather than the stated general behavior. The supplied background calls this a dated literature triage derived from OWR curation and explicitly requests independent verification of time-sensitive status and cited claims. No source visitation or independent literature assessment is asserted here.
\end{itemize}
\end{document}
